% Canonical node 020 · M002; current section path 1.3.5.
% Canonical owner: M002
\cnnaNodeHeading{020}{M002}{Birth-cut interior-domain theorem}
\cnnaNodePlacement{1.3.5}{D2}
Let $X_n$ be a recurrent C005 state, let $s_{n+1}$ be its C004 successor, and
let
\[
 \mathcal C_n(s_{n+1})=(B_n,I_n)
\]
be the canonical ordered M001 cut.  M001 fixes the coordinate sets and their
order, but it does not assign numerical entries to the four C006 blocks.
Consequently the cut alone cannot determine whether an interior matrix is
invertible or, equivalently for the C006 square interior system, whether
\[
 K_{II}X=K_{IB}
\]
has a unique exact fraction-value solution.

M002 therefore closes the second branch of its obligation: it states the exact
domain rather than inserting an unproved global invertibility hypothesis.

% CNNA-CONTEXT-BEGIN CTX-M002-PARTIAL-DOMAIN
This partiality is standard linear-algebraic hygiene: a Schur complement is not made total merely by notation.  Replacing the unique-solve domain with a pseudoinverse or regularization would define a different response operator and therefore a different model.

% CNNA-EXTREF-BEGIN EXT-USE-M002-PARTIAL-DIRECT
\cite[Basic Properties, pp.~17--46]{Zhang2005SchurComplement}.
% CNNA-EXTREF-END EXT-USE-M002-PARTIAL-DIRECT
% CNNA-CONTEXT-END CTX-M002-PARTIAL-DOMAIN
  For
explicit C006 block data $K$ whose dimensions agree with the M001 cut, define
\[
 \boxed{
 \mathcal D_{\mathrm{M002}}
 =\left\{(X_n,s_{n+1},K):
 \dim K=(|B_n|,|I_n|)\ \text{and}\
 K\in\operatorname{Dom}(\mathrm{C006})\right\}.}
\]
Here $K\in\operatorname{Dom}(\mathrm{C006})$ means exactly that the interior
system has one and only one encoded solve under the C006 fraction-value equation.  No floating tolerance,
condition-number threshold, pseudoinverse, regularization, positivity surrogate,
or matrix convention is substituted for this predicate.

The nonempty M001 boundary required by C006 is itself derived: the provenance
root is always on the M001 causal predecessor chain, hence $|B_n|>0$.
If $I_n=\varnothing$, C006 already gives a unique empty solve, so every block
assignment of the matching zero-interior shape belongs to the M002 domain.

For $|I_n|>0$, M001 alone cannot imply membership in the domain.  Indeed, for
the same fixed M001 dimensions choose all non-interior blocks zero and
$K_{IB}=0$.  With $K_{II}=I$ the unique solve is $X=0$; with $K_{II}=0$ every
matrix $X$ solves the homogeneous system, and because both $|B_n|$ and $|I_n|$
are positive there are at least two such solves.  Thus one and the same
combinatorial cut admits both an admissible and a non-admissible numerical block
assignment.  This obstruction is precisely why C007, not M001 or M002, must
construct the actual directed operator from the C005 conductances.

Python checks the exact M001 block dimensions and delegates acceptance exactly
to the C006 fraction-value admissibility predicate.  Lean enforces the same dimensions
in the type of \texttt{BirthCutBlocks}, proves the M001 boundary nonempty,
defines \texttt{InExactDomain} to be exactly C006 admissibility, proves the
zero-interior case, and transports domain membership to existence and uniqueness of the exact C006 response value modulo raw-fraction normalization.  The two languages therefore agree on the same exact
partial domain; the Python contrast witness is additional executable evidence
for the cut-only obstruction, not a different definition.

\paragraph*{Verified closure and limits.}
The theorem fixes the exact handoff rather than claiming universal
invertibility.  Boundary nonemptiness follows from the root port; zero
interior is unconditionally admissible; nonempty interiors require the stated
C006 unique-solve predicate.  The Python contrast witness demonstrates that
identical M001 dimensions can support both admissible and inadmissible block
entries.
